\section{Постановка задачи}
\subsection{Параметры ОУ и задающего воздействия}
Объект управления представлен в виде:
$$
\begin{cases}
    \dot{x}=Ax+bu \\
    y=Cx,
\end{cases}
$$
где $x\in\mathbb{R}^2$ - измеряемый вектор состояния, $u,y$ - измеряемые вход и выход объекта. Матрицы $A,b,C$ имеют вид:
$$
A=
\begin{bmatrix}
    0 & 3 \\
    4 & 7
\end{bmatrix},\quad
b=
\begin{bmatrix}
    0\\
    3
\end{bmatrix},\quad
C=
\begin{bmatrix}
    5 & 1
\end{bmatrix}
$$

\textbf{Цель:} синтез закона адаптивного управления, обеспечивающего ограниченность всех сигналов
и слежение выхода объекта за эталонным сигналом так, чтобы:
$$
\lim_{t\to\infty}\left(g\left(t\right)-y\left(t\right)\right)=0,
$$
где $g$ - мультисинусоидальное задающее воздействие с априори неизвестными амплитудами, частотами и
фазами гармоник. Для обеспечения ускоренной параметрической сходимости используется схема Лайона.

Параметры задающего воздействия $g(t)$:
$$
g(t)=9\cos\left(8t+1\right)\cos\left(6t-6\right)
$$

\subsection{Проверка объекта управления на свойство полной управляемости и 
наблюдаемости}
Для проверки ОУ на свойство полной управляемости воспользуемся критерием Калмана. Составим
матрицу управляемости:
$$
\mathcal{U}=
\begin{bmatrix}
    b & Ab
\end{bmatrix}=
\begin{bmatrix}
    0 & 9 \\
    3 & 21
\end{bmatrix},\quad
\operatorname{det}\mathcal{U}=0\cdot21-9\cdot3=-27\neq0
$$
Поскольку определитель квадратной матрицы $\mathcal{U}$ не нулевой, матрица имеет полный ранг
и система управляема.

Для проверки ОУ на свойство полной наблюдаемости воспользуемся критерием Калмана. Составим
матрицу наблюдаемости:
$$
\mathcal{O}=
\begin{bmatrix}
    C \\ CA
\end{bmatrix}=
\begin{bmatrix}
    5 & 1 \\
    4 & 22
\end{bmatrix},\quad
\operatorname{det}\mathcal{O}=5\cdot22-1\cdot4=106\neq0
$$
Поскольку определитель квадратной матрицы $\mathcal{O}$ не нулевой, матрица имеет полный ранг
и система наблюдаема.

\section{Перечень компонентов САУ и их параметры}
\section{Схема моделирования САУ с листингами расчётов}
\section{Переходные процессы САУ с градиентным АА}
\section{Переходные процессы САУ с модифицированным АА}
\section{Выводы}